
\documentclass[10pt]{article}
\usepackage[utf8]{inputenc}
\usepackage[T1]{fontenc}
\usepackage{amsmath, amssymb, xcolor, geometry, enumitem, multicol, tcolorbox}

\definecolor{mainblue}{RGB}{0, 102, 204}
\geometry{a4paper, margin=1.2cm, top=1cm, bottom=3cm, footskip=1.2cm}

\begin{document}
\enlargethispage{2.5cm}

% Modern Header
\begin{tcolorbox}[colback=mainblue!5, colframe=mainblue, sharp corners, boxrule=0.5pt]
    \small \textbf{Unit:} Unit 0: Foundations of Algebra \hfill \textbf{Name:} \rule{3cm}{0.4pt} \\
    \small \textbf{Topic:} Divisibility & Division Algorithm \hfill \textbf{Date:} \rule{2cm}{0.4pt} \\
    \centering \textbf{\large Homework (Jalapeno)}
\end{tcolorbox}

\vspace{0.2cm}

% MCQ Section
\begin{multicols}{2}
\begin{enumerate}[label=\textbf{Q\arabic*.}, leftmargin=*, itemsep=6pt, parsep=0pt]

  \item \small Is $24$ divisible by $3$?
  \begin{enumerate}[label=(\Alph*), leftmargin=0.6cm, nosep, itemsep=2pt]
    \item Yes
    \item No
    \item Maybe
    \item Depends on the context
    \item None of the above
  \end{enumerate}

  \item \small What is the remainder when $17$ is divided by $5$?
  \begin{enumerate}[label=(\Alph*), leftmargin=0.6cm, nosep, itemsep=2pt]
    \item $0$
    \item $1$
    \item $2$
    \item $3$
    \item $4$
  \end{enumerate}

  \item \small The ancient Egyptians built a pyramid with $180$ stone blocks in each row. How many rows can be made if each row must have a multiple of $10$ blocks?
  \begin{enumerate}[label=(\Alph*), leftmargin=0.6cm, nosep, itemsep=2pt]
    \item $10$ rows
    \item $15$ rows
    \item $18$ rows
    \item $20$ rows
    \item $25$ rows
  \end{enumerate}

  \item \small The Babylonians used a sexagesimal (base-$60$) number system. Is the number $120$ divisible by $4$ in this system?
  \begin{enumerate}[label=(\Alph*), leftmargin=0.6cm, nosep, itemsep=2pt]
    \item Yes
    \item No
    \item Maybe
    \item Depends on the context
    \item None of the above
  \end{enumerate}

  \item \small What is the result of $43 \div 5$ with remainder?
  \begin{enumerate}[label=(\Alph*), leftmargin=0.6cm, nosep, itemsep=2pt]
    \item $8$ with remainder $3$
    \item $9$ with remainder $1$
    \item $8$ with remainder $1$
    \item $7$ with remainder $2$
    \item $6$ with remainder $4$
  \end{enumerate}

  \item \small The ancient Greeks used $12$ gods to represent the $12$ months of the year. If $36$ stones are to be divided equally among these $12$ gods, how many stones will each god receive?
  \begin{enumerate}[label=(\Alph*), leftmargin=0.6cm, nosep, itemsep=2pt]
    \item $2$ stones
    \item $3$ stones
    \item $4$ stones
    \item $5$ stones
    \item $6$ stones
  \end{enumerate}

  \item \small A Mayan calendar has $365$ days in a year. If the days are divided into weeks of $7$ days, how many weeks are there in a year with remainder?
  \begin{enumerate}[label=(\Alph*), leftmargin=0.6cm, nosep, itemsep=2pt]
    \item $52$ weeks with remainder $1$ day
    \item $51$ weeks with remainder $2$ days
    \item $52$ weeks with remainder $3$ days
    \item $51$ weeks with remainder $1$ day
    \item $52$ weeks with remainder $2$ days
  \end{enumerate}

  \item \small The ancient Chinese used the remainder theorem to solve problems. If $17$ is divided by $5$, what is the remainder?
  \begin{enumerate}[label=(\Alph*), leftmargin=0.6cm, nosep, itemsep=2pt]
    \item $0$
    \item $1$
    \item $2$
    \item $3$
    \item $4$
  \end{enumerate}

\end{enumerate}
\end{multicols}

\vspace{-0.2cm}
\hrule
\vspace{0.2cm}
\textbf{\small Open Ended Questions (FRQ)}
\begin{enumerate}[label=\textbf{Q\arabic*.}, start=9, itemsep=5pt, topsep=0pt]

  \item \small The ancient Sumerians used a system of mathematics that included long division. Perform long division to find the quotient and remainder of $63 \div 9$. Show your work and write your answer as $\boxed{[quotient \text{ with remainder } remainder]}$.  \vspace{2cm}

  \item \small In ancient Egypt, a farmer has $48$ baskets of wheat to distribute among $8$ villages. If each village must receive an equal number of baskets, how many baskets will each village receive? Perform the division and write your answer as $\boxed{[answer]}$.  \vspace{2cm}

\end{enumerate}

\vfill

% Chef's Tips Footer
\begin{tcolorbox}[colback=blue!5, colframe=mainblue!50, title=\small \textbf{Chef's Tips}, fonttitle=\bfseries, bottom=1mm]
\begin{itemize}[noitemsep, topsep=2pt, leftmargin=0.5cm]
  \item \scriptsize Tip 1: Emphasize the importance of understanding the concepts of divisibility and division algorithm.\item \scriptsize Tip 2: Ensure students can apply the rules of divisibility for 2, 3, 5, 9, and 10.\item \scriptsize Tip 3: Encourage students to use real-world examples and historical contexts to practice problem-solving skills.
\end{itemize}
\end{tcolorbox}

\end{document}
  